%------------------------------------------------------------------------------------
%  ANNEXE C : GUIDE D'OUTILLAGE, COMPILATION ET BONNES PRATIQUES GIT
%------------------------------------------------------------------------------------
\chapter{Guide d'Outillage, Compilation et Bonnes Pratiques Git}
\label{annexe:guide_outillage}

La reproductibilité d'un travail scientifique impose une maîtrise complète des outils d'automatisation et de traçabilité.
Cette annexe détaille la chaîne de compilation et les conventions de versionnement recommandées pour le présent modèle.

\section{Le script d'automatisation PowerShell : \texttt{compile.ps1}}
\label{sec:script_compile}

Sous l'environnement Windows, les terminaux PowerShell utilisent historiquement une page de code OEM (CP 850 ou CP 1252) susceptible de dégrader l'affichage des caractères accentués dans les journaux de compilation.
Le script \texttt{scripts/compile.ps1} sécurise cette étape :

\begin{enumerate}
  \item Il sauvegarde la page de code d'entrée et de sortie de la console.
  \item Il bascule temporairement l'environnement en UTF-8 universel (Page de code 65001).
  \item Il déclenche l'utilitaire \texttt{latexmk} avec le moteur Lua\LaTeX{} et la redirection \texttt{-outdir=build}.
  \item Dès l'achèvement, il restaure la page de code d'origine de Windows sans altérer le reste de votre session.
\end{enumerate}

La commande d'exécution à la racine du projet est immédiate :
\begin{verbatim}
powershell -ExecutionPolicy Bypass -File scripts/compile.ps1
\end{verbatim}

\section{Nettoyage des fichiers intermédiaires : \texttt{clean.ps1}}
\label{sec:script_clean}

Durant les cycles de compilation, des dizaines de fichiers auxiliaires (\texttt{.aux}, \texttt{.bcf}, \texttt{.fls}, \texttt{.run.xml}, \texttt{.tdo}) s'accumulent dans le dossier \texttt{build/}.
Le script \texttt{scripts/clean.ps1} permet de purger ces résidus pour garantir une compilation vierge à froid sans supprimer vos images sources ni vos fichiers bibliographiques.

\section{Bonnes pratiques de versionnement sous Git}
\label{sec:bonnes_pratiques_git}

Un dépôt Git professionnel doit respecter des règles strictes de salubrité :
\begin{itemize}
  \item \textbf{Fichiers à versionner impérativement :}\\
  Tous les fichiers sources textuels (\texttt{*.tex}), la classe (\texttt{ClasseGuillaume.cls}),
  la base bibliographique (\texttt{references.bib}), les graphiques et images sources,
  les scripts PowerShell (\texttt{*.ps1}), et le fichier \texttt{.latexmkrc}.
  
  \item \textbf{Fichiers à ignorer formellement (\texttt{.gitignore}) :}
  La totalité du dossier \texttt{build/} et les extensions auxiliaires (\texttt{*.aux}, \texttt{*.log}, \texttt{*.toc}, \texttt{*.synctex.gz}, \texttt{*.pdf} intermédiaire).
  Seuls les fichiers sources doivent être conservés ; le document PDF final est un artefact dérivé qui se reconstruit de façon déterministe par simple compilation.

  \item \textbf{Diffs visuels avec \textsf{latexdiff} :}
  Grâce aux sauts de ligne sémantiques adoptés par ce modèle, vous pouvez utiliser l'utilitaire \texttt{latexdiff} pour générer automatiquement un document PDF de révision mettant en évidence en rouge les passages supprimés et en bleu les ajouts entre deux versions ou commits Git successifs.
\end{itemize}

